\section{Discussion}\label{sec:discussion}

Household simulation, rational or behaviorally extended, does not close the gap to the empirical benchmark. A structurally distinct hazard framework---no household utility, a substantially larger non-synthetic sample---recovers 91.3\% of it in microsimulation form on the benchmark-consistent accounting basis at the headline off-window floor and 97.9\% at the in-sample calibration point; across that floor's range recovery runs 88.2--93.7\%, a miss of 6.3 to 11.8 points. But no estimator's monthly path co-movement, Path B's included, survives removing the window's shared trend (Section~\ref{sec:pathb}), and the contrast is easy to over-read: the $\beta_1 = 0$ null recovers 85.7\% on the same basis at that headline floor, 88.7\% in-sample. Components common to any prepayment model---scheduled amortization, the seasoning baseline, the turnover floor---deliver the large majority of the level.

Named loan-level mechanisms are measured small wherever I can ablate them. Switching burnout off moves Path B by \$6.58 billion ($0.86$ points) at the in-sample floor and \$3.48 billion ($0.45$ points) at the headline off-window floor; zeroing the FICO/LTV covariates moves it $1.3$ points; the delinquency transition matrix is bounded at \$0.39 billion (Appendix~\ref{app:params}). The residual is modest at the in-sample calibration (2.1 points) but 8.7 points at the headline off-window calibration, roughly 1.6$\times$ the headline marginal. Three mechanisms remain \emph{candidates} for it, and for the timing failure every estimator shares, not demonstrated drivers of the level. Loan-level heterogeneity in credit score, geography, and life circumstance makes mobility no smooth function of coupon and market rate alone; the hazard framework's 295 stratum fixed effects let Path A track it where the ABM's coupon-cohort buckets cannot. Servicer behavior---loss mitigation, forbearance, and modification, unusually active in the higher-rate, higher-delinquency-risk QT window---reaches prepayment and default outcomes outside any borrower-side decision model; Path B's competing-risks structure and Markov delinquency pipeline represent it only partially, and nothing beyond that pipeline is modeled at all. Burnout is representable at the cohort level, as in both hazard paths, but not inside a closed, per-agent household population (Section~\ref{sec:abm-results}'s falsified survivor-selection experiment). Servicer behavior beyond the pipeline, updated-equity dynamics, and home-equity extraction carry no bound of any kind here, each plausibly rate-cycle-correlated.

\subsection{The Dual Systemic Fallout}\label{sec:discussion-fallout}

Transmission into the fixed-income sector creates a two-pronged systemic risk. The first channel runs through the central bank itself: extending MBS durations deprives the SOMA of its scheduled cash flows and breaks the predictability effective Quantitative Tightening requires. The \$764.7 billion shortfall-versus-cap is mostly rate-inelastic; the rate-driven object is the duration extension itself, a 9.4-year portfolio WAL under the empirical path against 3.4 years at 2021 realized speeds (Appendix~\ref{sec:robustness-wal}). That 3.4-year comparator is a refinancing-boom baseline. The extension, not the cap arithmetic, feeds the fragility channel below.

The second channel runs through private banks, which, unlike the Fed, must mark to market: the structural duration extension of these low-coupon MBS is one source of the unrecognized mark-to-market losses accumulated across the banking sector. \citet{jiang2023} show duration extension alone was not sufficient to precipitate failure---it was mark-to-market losses interacting with a high share of uninsured deposits that made specific regional institutions vulnerable to self-fulfilling runs. The extension, a necessary precondition for that channel, comes from the below-market coupon stock meeting rate normalization under any payoff rule: Appendix~\ref{sec:abm-danish}'s rule-only counterfactual shortens the portfolio by only 0.6 of the roughly six extension-years (Appendix~\ref{sec:robustness-wal}), and those six years are measured against 2021's refinancing-boom speeds, not normal turnover. Read instead against the baseline turnover floor, the empirical path's approximate weighted-average life is 9.4 years against 9.5 at the window open and 8.5 against 8.6 at its close, so the extension measures refinancing-versus-turnover, not lock-in-versus-normalcy. The elasticity channel I identify is therefore a second-order contributor to that precondition, and neither it nor the extension is, on its own, a sufficient condition for bank failure. I price cash-flow timing alone (principal arriving later than the cap schedule assumed), not a market-value loss: Appendix~\ref{sec:robustness-wal} declines option-adjusted duration and convexity as out of scope, so no mark-to-market or DV01 image is put on that 0.6 of an extension-year, the bank-side losses \citet{jiang2023} price have no counterpart here, and the fragility channel is invoked as context, not as something this design measures.

The lock-in effect also imposes a real economic cost by keeping households from relocating toward higher-productivity labor markets. \citet{hsieh2019} estimate substantial aggregate output losses from such spatial misallocation of labor under housing supply constraints in high-productivity U.S. cities; their headline magnitudes are cumulative 1964--2009 level effects revised in a published correction, so I cite the mechanism, not a point estimate. Theirs is a supply-side housing-availability constraint; the demand-side mobility constraint I quantify runs through the same channel, whatever the housing supply elsewhere. No per-household or per-move welfare estimate enters the paper---the per-move surplus that would turn this leg into a magnitude is nowhere quantified---so the labor-reallocation leg is directional only. Nor do I translate the \$764.7 billion shortfall into welfare terms: it is an accounting gap against a cap schedule, mostly rate-inelastic by the decomposition above, and neither a lower nor an upper bound on the par-payoff rule's welfare cost. The welfare-relevant objects are the constrained moves, whose fiscal image is the lock-in marginal ($+\$42.6$ billion of trapped roll-off at the headline off-window floor, $+\$70.3$ billion at the in-sample calibration point), and the forgone labor-reallocation gains this literature identifies, which my accounting framework does not measure.

A back-of-envelope prices the extension in fiscal units. At the window's policy-rate levels the funding spread over the book's 2.49\% weighted-average coupon ran roughly 250--300 basis points, so the \$42.6 billion of trapped roll-off carries on the order of \$1 billion per year of negative carry on the consolidated public balance sheet, for as long as the gap persists. That restates in remittances terms the maturity extension Table~\ref{tab:wal} prices in years; it derives from committed quantities---the book coupon, the marginal---and the public policy-rate path, and is subject to the same face-versus-cash incidence reading as the buyback bracket (Appendix~\ref{app:params}).

\subsection{Evaluating Alternative Securitization Policies}\label{sec:discussion-policy}

Dividing the central-minus-null differential by the mean balance of the sample's surviving loans (\$237{,}317 across 40{,}077 loans) gives about 179{,}500 foregone payoffs on the SOMA book over the 42 window months, or 51{,}300 a year, both at the production $\mu = 1$ convention. That convention treats the whole voluntary gap response as the moving channel and therefore returns an \emph{upper bound}; the all-loan mean averages in prepaid loans at zero balance and would roughly double the count. Confining the response to a half or a quarter of the voluntary hazard gives 108{,}300 and 58{,}700 (run \texttt{marginal\_transaction\_counts}). Grossing to the whole market on a SOMA share of 20--22\% of 1--4 family mortgage debt puts the upper-bound count at 233{,}000--256{,}000 payoffs a year, or 5.7--6.3\% of the 4.08 million annual existing-home-sale run rate. This is the first outcome-side check the design admits, and even at its upper bound the identified channel is a small minority of transaction volume. \citet{fonseca2026} document a 40\% fall in existing home sales between 2022 and 2024, but that series has no pre-2025 history on this project's data access, so the marginal's share of \emph{that decline} cannot be formed without assuming the 2024 level equals the 2025 one. I decline the assumption and report the run-rate share instead.

In response to the mobility collapse, U.S. policymakers revisited assumable and portable mortgages. In late 2025 the Federal Housing Finance Agency stated that it was actively evaluating portability, and that Fannie Mae and Freddie Mac were examining wider assumption; both permit fixed-rate assumption only in limited circumstances such as death or divorce, and no rulemaking had issued as of writing. I therefore cite this as the policy setting, not an enacted rule, and what follows as a taxonomy of that design space, not a finding; Section~\ref{sec:intro}'s concession applies to any evaluation of it. The statutory right already exists for the 20.4\% of SOMA face that is Ginnie Mae collateral, but realized take-up is measured nowhere in this design and in no public assumption-volume series I could find, so the 20.4\% share is used throughout as a ceiling on the carve-out, not an estimate of it. Both instruments explicitly prevent prepayment upon relocation: under assumption the below-market coupon transfers to the buyer, who must separately finance the gap between purchase price and assumed balance, typically at prevailing rates, which limits take-up. They therefore relocate duration extension risk, solving the microeconomic mobility constraint while cementing the macroeconomic extension risk onto the investor's balance sheet.
